% chapters/conclusion.tex
\chapter{Conclusion and Future Work}
\label{ch:conclusion}

\section{Summary of Contributions}
\label{sec:conc-summary}

This dissertation introduced a three-stage human-centred framework for evaluating
causal explanations produced by Large Language Models. The framework combines
AHP-derived criterion weights with adaptive question selection and interval-based
aggregation of Likert ratings, producing a structured and reproducible method for
comparing LLMs across multiple explanation quality dimensions.

Applied to nine open-source LLMs evaluated by graduate student reviewers on economic
reasoning tasks, the framework produced four main findings.

First, graduate students prioritise trustworthiness and completeness when evaluating
LLM explanations. Together, these two criteria account for more than 75\% of total
criterion weight, a pattern that standard accuracy benchmarks do not capture.

Second, interval-based aggregation reveals reasoning instability that scalar means
conceal. The performance floor of the top model (\texttt{qwen2.5:14b}) far exceeds
that of the weakest model (\texttt{codellama:7b}), and interval width is highest in
the task categories where distinguishing valid from plausible reasoning is hardest.

Third, a systematic fluency trap is identified: several models achieve relatively
high understandability scores while receiving much lower trustworthiness scores.
This pattern is most pronounced in \texttt{mistral:7b} and carries direct
implications for educational deployment, where students may misplace trust in fluent
but causally unreliable explanations.

Fourth, no single model excels uniformly across all five reasoning task categories,
indicating that model selection should be task-specific rather than based on
aggregate rankings.

\section{Broader Impact}
\label{sec:conc-impact}

The broader impact of this work lies in providing researchers and practitioners with
a principled process for selecting LLMs and for designing explanations that better
serve student needs. The combination of human ratings, dispersion-based sampling,
and interval-based aggregation can support more transparent and defensible model
selection in settings where explanation quality has direct consequences for learning
outcomes or decision-making.

The identification of the fluency trap is particularly consequential for AI-assisted
education. Deploying models that produce coherent-sounding but causally unreliable
explanations risks reinforcing student misconceptions rather than correcting them.
Evaluation frameworks that measure trustworthiness explicitly — rather than relying
on accuracy or fluency as proxies — are therefore essential for responsible
deployment in academic contexts.

\section{Future Work}
\label{sec:conc-future}

Several directions extend the present framework.

\textbf{Broader model and domain coverage.} Applying the framework to closed-source
models and to disciplines beyond economics would test the generalisability of the
criterion weights and the ranking procedure. Different academic domains may produce
different weight profiles, reflecting different conventions for acceptable causal
argumentation.

\textbf{Adaptive weighting across reviewer populations.} The current AHP weights are
derived from a single cohort of graduate students. Future work could investigate
whether weights differ systematically across student populations (e.g., undergraduate
versus doctoral, different disciplinary backgrounds) and develop adaptive weighting
strategies that reflect these differences.

\textbf{Longitudinal evaluation.} The current study evaluates explanation quality at
a single point in time. Longitudinal studies that track how student ratings of the
same model evolve over an academic term could reveal whether repeated exposure to
LLM explanations affects trust calibration or criterion weighting.

\textbf{Explanation-aware fine-tuning.} The framework's outputs — human ratings
across structured criteria — could be used as feedback signals for explanation-aware
fine-tuning. Models trained to optimise for trustworthiness and completeness, as
measured by student reviewers, may produce explanations that are more aligned with
educational needs than models optimised solely for accuracy.

\textbf{Automated quality monitoring.} Building on the dispersion metrics developed
in Stage B, automated tools could monitor explanation quality in deployment by
flagging outputs where structural dispersion across repeated queries is high,
signalling potential causal inconsistency without requiring human review of every
instance.
